\chapter{Implementação}
\label{cap:implementacao}

\orientacao{Este capítulo apresenta a materialização da solução projetada em software, descrevendo o ambiente utilizado, a organização da aplicação, a implementação dos componentes, sua integração e os procedimentos de teste. Sua função é demonstrar como o projeto definido nos capítulos anteriores foi efetivamente transformado em uma solução executável e quais resultados foram obtidos.}

\section{Ambiente de Desenvolvimento}
\orientacao{Descrever o ambiente utilizado para desenvolver e executar o sistema, incluindo hardware, sistema operacional, ferramentas e configurações relevantes.}

\section{Estrutura da Aplicação}
\orientacao{Apresentar a organização dos arquivos, módulos, pacotes e demais elementos que compõem a aplicação implementada.}

\diagrama{Diagrama de pacotes}{pacotes}

\section{Implementação dos Componentes}
\orientacao{Descrever como os componentes definidos no projeto foram efetivamente implementados, destacando aspectos técnicos relevantes.}

\section{Integração dos Módulos}
\orientacao{Descrever como os diferentes módulos e componentes da aplicação foram integrados e como ocorre sua comunicação em ambiente de execução.}

\diagrama{Diagrama de implantação}{implantacao}

\section{Testes Realizados}
\orientacao{Apresentar os procedimentos de teste realizados, os aspectos avaliados e os principais resultados obtidos.}

\section{Resultados Obtidos}
\orientacao{Apresentar os resultados alcançados com a implementação e verificar em que medida eles atendem aos objetivos e às necessidades identificadas.}

\section{Limitações da Implementação}
\orientacao{Registrar limitações, funcionalidades não implementadas, problemas conhecidos e aspectos que possam exigir complementação ou evolução futura.}
